% Vetores unitários em coordenadas cilíndricas
% Feito para o ÉguaTeX (eguatex.com.br). Uso livre (CC0 1.0).
%
% Compile e pronto: sai só a figura, em PDF.
% Para pôr no seu trabalho, use o PDF com \includegraphics ou copie
% o bloco tikzpicture para dentro de uma figure.
%
% Os vetores mudam de direção de ponto para ponto: dependem de phi.
\documentclass[tikz,border=4pt]{standalone}
\usepackage[T1]{fontenc}
\usepackage{lmodern}
\usepackage{amsmath}

\colorlet{corVet}{red!75!black}

\begin{document}
\begin{tikzpicture}[
  x={(-0.4435cm,-0.3663cm)}, y={(0.9768cm,-0.1715cm)}, z={(0cm,1cm)},
  line cap=round, line join=round, >=stealth,
  eixo/.style={line width=0.75pt},
  guia/.style={line width=0.55pt, black!35, dashed},
  vetor/.style={->, line width=1.4pt, corVet}]

% Mude estes números e o desenho e os rótulos acompanham.
\def\praio{3.2} \def\pang{50} \def\pz{2.4} % ponto P = (rho, phi, z), phi em graus
\def\tam{1.3}                           % tamanho dos vetores, em cm
\pgfmathsetmacro{\c}{cos(\pang)} \pgfmathsetmacro{\s}{sin(\pang)}

% eixos
\draw[eixo] (0,0,0) -- (5,0,0) node[below left] {$x$};
\draw[eixo] (0,0,0) -- (0,5.5,0) node[right] {$y$};
\draw[eixo] (0,0,0) -- (0,0,4.8) node[above] {$z$};

% guias até o ponto
\draw[guia] (0,0,0) -- ({\praio*\c},{\praio*\s},0) -- ({\praio*\c},{\praio*\s},\pz);
\draw[guia] (0,0,\pz) -- ({\praio*\c},{\praio*\s},\pz);
\draw[guia] plot[domain=0:\pang, samples=30] ({\praio*cos(\x)},{\praio*sin(\x)},0);
\draw[->] plot[domain=0:\pang, samples=30] ({1.3*cos(\x)},{1.3*sin(\x)},0);
\node at ({1.7*cos(\pang/2)},{1.7*sin(\pang/2)},0) {$\phi$};
\path (0,0,\pz) -- ({\praio*\c},{\praio*\s},\pz) node[midway, above] {$\rho$};
\fill ({\praio*\c},{\praio*\s},\pz) circle[radius=1.8pt] node[left=3pt] {$P$};

% vetores unitários no ponto P
\coordinate (P) at ({\praio*\c},{\praio*\s},\pz);
\draw[vetor] (P) -- ++({\tam*\c},{\tam*\s},0) node[below=1pt] {$\hat{\boldsymbol{\rho}}$};
\draw[vetor] (P) -- ++({-\tam*\s},{\tam*\c},0) node[right] {$\hat{\boldsymbol{\phi}}$};
\draw[vetor] (P) -- ++(0,0,\tam) node[above] {$\hat{\mathbf{z}}$};
\end{tikzpicture}
\end{document}
